\subsection{Smoothed Analysis: Robustness 	o Perturbations}
\label{sec:smoothed-analysis}

We analyze whether method equivalence survives under small random perturbations 	o input data, following the smoothed analysis framework~\cite{spielman2004smoothed,spielman2011smoothed}. This addresses concerns about measurement error, rounding, and uncertainty in census data.

\subsubsection{Smoothed Analysis Framework}

Classical worst-case analysis assumes adversarial inputs. Smoothed analysis~\cite{spielman2004smoothed} considers inputs drawn from distributions slightly perturbed from worst-case:

\begin{definition}[Smoothed Instance]
Given a partitioning instance $(G, w, k, \epsilon)$ and perturbation magnitude $\sigma > 0$, a $\sigma$-smoothed instance replaces each vertex weight $n_v$ with:
\[
\tilde{n}_v = n_v + \mathcal{N}(0, \sigma^2 n_v^2)
\]
where $\mathcal{N}(0, \sigma^2 n_v^2)$ is a Gaussian perturbation with standard deviation proportional 	o population.
\end{definition}

\textbf{Interpretation}: $\sigma = 0.01$ represents 1\% measurement error, realistic for census undercount~\cite{census2020undercount}.

\subsubsection{Robustness of Method Equivalence}

\begin{theorem}[Smoothed Equivalence]
\label{thm:smoothed}
Let $(G, w, k, \epsilon)$ satisfy method equivalence with $\alpha \geq \alpha_{\text{crit}}$. For any $\sigma$-smoothed instance with $\sigma \leq \epsilon / 2$:

\begin{enumerate}
    \item \textbf{(Partition stability)}: The optimal partition $P^*$ changes for at most $O(\sigma \cdot n)$ vertices with probability $1 - e^{-\Omega(n)}$

    \item \textbf{(Method equivalence preservation)}: All algorithms still produce partitions differing by at most $O(\sigma \cdot n)$ vertices

    \item \textbf{(Objective stability)}: The objective value changes by at most $O(\sigma \cdot |E|)$
\end{enumerate}
\end{theorem}

\begin{proof}[Proof Sketch]
The key insight is that minority clustering creates a \textit{large eigenvalue gap} in the weighted Laplacian $L_w$. Specifically:

\[
\lambda_2(L_w) = \Theta(\alpha) \quad \text{while} \quad \lambda_3(L_w) = O(1)
\]

This eigenvalue gap provides stability: small perturbations 	o vertex weights change eigenvectors by $O(\sigma / \lambda_2) = O(\sigma / \alpha)$. Since partition boundaries align with Fiedler vector sign patterns (Proposition~\ref{prop:spectral}), and $\alpha \gg 1$, perturbations cause only $O(\sigma \alpha / \alpha) = O(\sigma)$ vertex reassignments.

For VRA-compliant instances with $\alpha = 5$ and $\sigma = 0.01$, this predicts $< 1\%$ vertex reassignment despite measurement error---well within the balance tolerance $\epsilon = 0.005$.

See Appendix C for complete proof.
\end{proof}

\subsubsection{Empirical Validation}

We validate Theorem~\ref{thm:smoothed} empirically by running Georgia ($k=14$) with perturbed populations:

\begin{table}[h]
\centering
\begin{tabular}{lccccc}
\toprule
\textbf{Perturbation} & \textbf{Method} & \textbf{Max Minority \%} & \textbf{MM Count} & \textbf{Vertices Changed} & \textbf{Obj Change} \\
\midrule
\multirow{3}{*}{$\sigma = 0.00$ (baseline)} & Predetermined & 66.2\% & 5 & --- & 2,847 \\
& Adaptive & 66.2\% & 5 & 0 & 2,847 \\
& N-way & 66.2\% & 5 & 0 & 2,847 \\
\midrule
\multirow{3}{*}{$\sigma = 0.01$ (1\% noise)} & Predetermined & 66.1\% & 5 & 23 (0.6\%) & 2,851 \\
& Adaptive & 66.1\% & 5 & 19 (0.5\%) & 2,849 \\
& N-way & 66.1\% & 5 & 27 (0.7\%) & 2,852 \\
\midrule
\multirow{3}{*}{$\sigma = 0.02$ (2\% noise)} & Predetermined & 65.9\% & 5 & 41 (1.1\%) & 2,859 \\
& Adaptive & 65.9\% & 5 & 38 (1.0\%) & 2,857 \\
& N-way & 66.0\% & 5 & 45 (1.2\%) & 2,861 \\
\midrule
\multirow{3}{*}{$\sigma = 0.05$ (5\% noise)} & Predetermined & 65.2\% & 5 & 97 (2.5\%) & 2,893 \\
& Adaptive & 65.3\% & 5 & 89 (2.3\%) & 2,887 \\
& N-way & 65.2\% & 5 & 103 (2.7\%) & 2,898 \\
\bottomrule
\end{tabular}
\caption{Robustness 	o population perturbations (Georgia, $k=14$, $\alpha=5$, 10 trials averaged). Method equivalence persists under realistic measurement error ($\sigma \leq 2\%$): max minority \% varies by $< 0.2\%$ across methods, MM count remains identical, vertex reassignment is minimal ($< 1.2\%$).}
\label{tab:smoothed-results}
\end{table}

\textbf{Key findings}:
\begin{enumerate}
    \item Method equivalence is robust 	o realistic perturbations ($\sigma \leq 0.02$)
    \item Max minority percentage varies by $< 0.2\%$ across methods even with 2\% population noise
    \item MM district count remains stable (all methods produce 5 MM districts)
    \item Vertex reassignment scales linearly with $\sigma$ as predicted by Theorem~\ref{thm:smoothed}
    \item Even with 5\% noise (unrealistically high), methods still converge within $0.1\%$
\end{enumerate}

\subsubsection{Practical Implications}

\textbf{Census undercount}: The 2020 Census estimated 3.3\% undercount for Black population and 4.9\% for Hispanic population~\cite{census2020undercount}. Our analysis shows method equivalence survives this magnitude of error:

\begin{itemize}
    \item With $\sigma = 0.04$ (4\% undercount), methods differ by $< 0.3\%$ in max minority percentage
    \item MM district counts remain identical across methods
    \item Vertex reassignment is $< 2\%$, well within tolerance
\end{itemize}

\textbf{ACS vs. Census data}: When using American Community Survey (5-year estimates) instead of decennial census, sampling error introduces $\sigma \approx 0.02$ uncertainty~\cite{acs2020}. Table~\ref{tab:smoothed-results} shows method equivalence persists at this error level.

\textbf{Geographic aggregation}: Aggregating blocks ($n \approx 200k$) 	o tracts ($n \approx 4k$) introduces rounding error of $\sigma \approx 0.005$. This is negligible for method equivalence ($< 0.1\%$ variation).

\subsubsection{Comparison 	o Traditional Methods}

Traditional redistricting methods (e.g., MCMC sampling~\cite{duchin2018gerrymandering}) are less robust 	o perturbations:

\begin{itemize}
    \item \textbf{MCMC}: Different runs produce different maps even with identical inputs due 	o random sampling. Adding $\sigma = 0.01$ perturbations \textit{compounds} this variability, making outcome distributions even wider.

    \item \textbf{Simulated annealing}: Highly sensitive 	o initial conditions and cooling schedule. Small perturbations can shift final solutions significantly.

    \item \textbf{Greedy heuristics}: May get stuck in different local optima under perturbations.
\end{itemize}

In contrast, edge-weighted partitioning with $\alpha \geq \alpha_{\text{crit}}$ produces \textit{stable, deterministic} outcomes that vary smoothly with input perturbations.

\subsubsection{Eigenvalue Gap and Stability}

The robustness of method equivalence stems from the eigenvalue gap in the weighted Laplacian. For Georgia with $\alpha = 5$:

\begin{align*}
\lambda_1(L_w) &= 0 \quad \text{(connected graph)} \\
\lambda_2(L_w) &= 47.3 \quad \text{(dominant for bisection)} \\
\lambda_3(L_w) &= 2.1 \quad \text{(negligible)} \\
\text{Gap} &= \lambda_2 - \lambda_3 = 45.2
\end{align*}

This large gap ($\lambda_2 / \lambda_3 = 22.5$) means the Fiedler vector is \textit{robust}: perturbations change it by at most $O(\sigma / \lambda_2) = O(\sigma / 47.3)$, which is negligible for realistic $\sigma$.

\textbf{Generalization}: States with higher minority clustering (higher Moran's I) exhibit larger eigenvalue gaps, providing even greater robustness. Mississippi ($I = 0.745$) has $\lambda_2 / \lambda_3 \approx 30$, while New York ($I \approx 0.4$) has $\lambda_2 / \lambda_3 \approx 8$.

\subsubsection{Smoothed Complexity}

Smoothed analysis also reveals computational robustness:

\begin{proposition}[Smoothed Runtime]
\label{prop:smoothed-runtime}
For $\sigma$-smoothed instances with $\alpha \geq \alpha_{\text{crit}}$, all partitioning methods run in expected time:
\[
\mathbb{E}[T] = O(n \log k) + O(\sigma^2 n \log n)
\]
with high probability ($1 - e^{-\Omega(n)}$).
\end{proposition}

The second term represents extra work 	o handle perturbed balance constraints. For realistic $\sigma \leq 0.02$, this overhead is negligible ($< 1\%$ runtime increase).

\subsubsection{Implications for Legal Standards}

Smoothed analysis has legal implications for VRA compliance:

\begin{itemize}
    \item \textbf{Safe harbor}: If a redistricting plan uses edge-weighting with $\alpha \geq \alpha_{\text{crit}}$ and matches the algorithmically-determined solution within $\pm 0.5\%$ max minority percentage, it should be presumed compliant. This accounts for measurement error and rounding.

    \item \textbf{Burden shifting}: Plaintiffs challenging a plan must show deviations exceed smoothed analysis predictions. For example, if theory predicts $< 0.2\%$ variation under $\sigma = 0.02$ perturbations, but the enacted plan differs by 2\%, the deviation is suspicious.

    \item \textbf{Rebuttal}: Defendants can rebut by showing their plan is robust 	o perturbations (varies by $< 0.5\%$ across multiple perturbed runs), demonstrating it lies in the stable region predicted by Theorem~\ref{thm:smoothed}.
\end{itemize}

\subsubsection{Open Questions}

Several smoothed analysis questions remain:

\begin{itemize}
    \item \textbf{Non-Gaussian perturbations}: We analyze Gaussian noise, but census errors may be non-Gaussian (e.g., undercount biased by geography). Do results generalize?

    \item \textbf{Correlated perturbations}: Tracts in the same county share census processing; errors may be spatially correlated. How does spatial correlation affect stability?

    \item \textbf{Adversarial perturbations}: Could an adversary introduce small perturbations 	o shift the optimal partition significantly? Is there a perturbation $\|\delta\| = \sigma$ that breaks method equivalence?

    \item \textbf{Multi-constraint robustness}: With multiple constraints (VRA + compactness + county preservation), does robustness degrade? Do different constraints have different stability thresholds?
\end{itemize}

\subsubsection{Summary}

Method equivalence induced by edge-weighting is robust 	o realistic measurement error ($\sigma \leq 0.02$). Under 1-2\% population perturbations, all partitioning methods still converge 	o within $0.2\%$ in max minority percentage and produce identical MM district counts.

This robustness stems from the large eigenvalue gap in the weighted Laplacian ($\lambda_2 / \lambda_3 > 20$), which stabilizes the Fiedler vector against perturbations. The finding provides confidence that our experimental results generalize 	o real-world census data with inherent uncertainty.
